\documentclass{article}
\usepackage{graphicx} % Required for inserting images
\usepackage{hyperref}

\title{Quantifying Forest Loss from Wildfires - Case Study}
\author{Jinal P, Gianluca F.}
\date{October 2026}

\begin{document}

\maketitle

\section{Introduction}

\subsection{The spread of wildifires}
On July 12, 2022, a thermal vehicle van parked outside the pine forest in La Teste-de-Buch broke down and caught on fire. The unfortunate location coupled with a persistent heatwave across Europe, a drought and lowest amount of precipitation the region experienced since 1959  caused the fire to gain quick speed and spread across the forest.  A second fire broke out in Landiras as well. Around 2,000 firefighters combined with Airforce Military units, help from neighbouring nations and local volunteers managed to curb the fire by its 11th day in La Teste-de-Buch and 13th day in Landiras. The fire reappeared in August in Landiras through an underground spread due to peat deposits and took further effort to be stopped. 

\subsection{Aftermath}
The authorities reported a total of 20,800 hectares of forest burned in fire, with 13,800 and 7,000 hectares in Landiras and La Teste-de-Buch respectively.  Around 36,750 people were evacuated with loss of 3 homes in La Teste-de-Buch and 2 homes in Landiras. While no human life was lost to the fire, the region suffered a huge loss in terms of flora, fauna and local animal life.

\subsection{Importance of Remote Sensing}
As extreme weather events such as droughts and heatwaves become more frequent due to climate change, the use satellite imagery to support detection and monitoring of wildfires to coordinate response efforts is crucial moving forward. It holds an important role in aftermath analysis as well, as exemplified through satellite imagery used on Wikipedia Page to showcase the spread and reach of the wildfire.


\end{document}
